\section{Original shells and marked words}

Fix a prime
\[
 p\equiv1\pmod{24}
\]
and an integer $a$ with $p/4<a<p/2$.  Retain
\[
 R=4a-p,\qquad S=pa.
\]
Then $0<R<p$, $R\equiv3\pmod4$, and
\[
 \gcd(p,R)=\gcd(a,R)=\gcd(S,R)=1.                 \tag{1.1}
\]
The last equalities follow from $R=4a-p$, $a<p$, and the primality of $p$.
Every divisor of $S$ therefore represents a unit modulo $R$.

For $u\mid a^2$, put
\[
 g_u=\gcd(a,u),\qquad
 (h_u,r_u,s_u)=
 \left(\frac{g_u^2}{u},\frac{u}{g_u},\frac{a}{g_u}\right).       \tag{1.2}
\]
These are positive integers and retain the exact inverse data
\[
 a=h_ur_us_u,\qquad u=h_ur_u^2,\qquad \gcd(r_u,s_u)=1.            \tag{1.3}
\]
The original full word sets are
\[
 E_a=\{u\mid a^2:R\mid4u+1\},\qquad
 M_a=\{u\mid a^2:R\mid p+4u\}.                                 \tag{1.4}
\]
Their marked ordered returns are
\begin{align}
 E_a:\quad&\kappa=\frac{pr+s}{R},&
 (x,y,z)&=(a,hs\kappa,phr\kappa),                               \tag{1.5}\label{eq:Ereturn}\\
 M_a:\quad&\lambda=\frac{r+s}{R},&
 (x,y,z)&=(a,phs\lambda,phr\lambda).                            \tag{1.6}\label{eq:Mreturn}
\end{align}
The classical type-I/type-II parametrizations are the literature antecedent
\citep[Section~2, Propositions~2.2 and~2.6]{ElsholtzTao}.  Equations
\eqref{eq:Ereturn}--\eqref{eq:Mreturn} retain the marked coordinates used in
this work; no existence statement is imported from the parametrization.

\section{The complete original-source signed pairing}

Let $G_R=(\mathbb Z/R\mathbb Z)^\times$.  In the real Hilbert space with
orthonormal basis $\{e_t:t\in G_R\}$, define, for every real $\sigma$,
\[
 v_{p,a,\sigma}
 =\sum_{d\mid S}\left(\frac d{\sqrt S}\right)^{-\sigma}e_{d\bmod R},
 \qquad J_Re_t=e_{-t}.                                          \tag{2.1}
\]
Different divisors with the same residue remain as separate summands before
their coefficients are added at the same basis vector.  The vector exists at
every shell, whether or not either gate in (1.4) is occupied.  Set
\[
 \mathcal T_{p,a}(\sigma)
 =\langle v_{p,a,\sigma},J_Rv_{p,a,\sigma}\rangle
 =\sum_{\substack{d,e\mid pa\\R\mid d+e}}
   \left(\frac{de}{pa}\right)^{-\sigma}.                         \tag{2.2}\label{eq:signedpair}
\]
The unweighted pair statistic and its character form occur in the earlier
Boolean-product collaboration \citep[BP15--BP16]{Boolean}; the next theorem
computes all weighted fibres over the original marked words.

\begin{theorem}[Exact positive-weight fibres]\label{thm:weighted-fibres}
For
\[
 W_\sigma(h)=\sum_{k\mid h}\left(\frac h{k^2}\right)^\sigma,
\]
one has, for every real $\sigma$,
\[
 \boxed{
 \mathcal T_{p,a}(\sigma)
 =2\sum_{u\in E_a}W_\sigma(h_u)
 +(p^\sigma+p^{-\sigma})\sum_{u\in M_a}W_\sigma(h_u).}          \tag{2.3}\label{eq:fibreformula}
\]
Consequently
\[
 \boxed{\mathcal T_{p,a}(\sigma)>0
 \Longleftrightarrow E_a\cup M_a\ne\varnothing.}               \tag{2.4}
\]
At $\sigma=0$ this becomes
\[
 \mathcal T_{p,a}(0)
 =2\sum_{u\in E_a}\tau_{\rm div}(h_u)
 +2\sum_{u\in M_a}\tau_{\rm div}(h_u).                         \tag{2.5}
\]
\end{theorem}

\begin{proof}
Because $p\nmid a$, every divisor of $pa$ is uniquely $p^\epsilon b$ with
$\epsilon\in\{0,1\}$ and $b\mid a$.  For an ordered pair $b,c\mid a$, retain
\[
 k=\gcd(b,c),\qquad b=kr,\qquad c=ks,\qquad\gcd(r,s)=1.
\]
The least common multiple $krs$ divides $a$.  Writing $h=a/(rs)$, the exact
remaining condition is $k\mid h$.  Conversely every $k\mid h$ returns the
actual divisors $kr,ks\mid hrs=a$.  The associated original marked divisor is
\[
 u=hr^2,\qquad a=hrs.                                            \tag{2.6}
\]

The patterns with neither or both entries divisible by $p$ are
\[
 (kr,ks),\qquad(pkr,pks),\qquad k\mid h.
\]
Because $k$ is a unit modulo $R$, their opposition congruence is $R\mid r+s$.
Using $p=4hrs-R$ and the fact that $4hr$ is a unit modulo $R$, this is equivalent
to $R\mid p+4hr^2$, the middle gate.  The two retained weights are
\[
 p^\sigma\left(\frac h{k^2}\right)^\sigma,\qquad
 p^{-\sigma}\left(\frac h{k^2}\right)^\sigma.                  \tag{2.7}
\]

The mixed patterns are
\[
 (pkr,ks),\qquad(ks,pkr),\qquad k\mid h.
\]
Here $r$ is always the coprime part carried by the $p$-bearing divisor; the
second displayed pair retains the opposite position bit.  Conversely, for any
mixed ordered pair, remove the one factor $p$, orient $r$ from that entry and
$s$ from the other entry, and then recover $k,h,u$.  This gives both position
bits without omitting or duplicating an ordered pair.
Their opposition congruence is $R\mid pr+s$.  Since
\[
 pr+s=s(4hr^2+1)-Rr
\]
and $s$ is a unit modulo $R$, this is equivalent to the exterior gate
$R\mid4u+1$.  Each retained weight is
\[
 \left(\frac h{k^2}\right)^\sigma.                              \tag{2.8}
\]
The four $p$-patterns exhaust every ordered pair in (2.2), proving
\eqref{eq:fibreformula}.  Every summand is strictly positive for real $\sigma$,
which proves (2.4) and (2.5).
\end{proof}

\begin{proposition}[Character sign, parity energies, and degree action]
For a character $\chi\in\widehat G_R$, retain
\[
 Z_{S,\chi}(\sigma)=
 \sum_{d\mid S}\chi(d)\left(\frac d{\sqrt S}\right)^{-\sigma}.
\]
Then
\[
 \boxed{\mathcal T_{p,a}(\sigma)=
 \frac1{\varphi(R)}\sum_{\chi\in\widehat G_R}
 \chi(-1)|Z_{S,\chi}(\sigma)|^2.}                               \tag{2.9}\label{eq:char}
\]
With $P_\pm=(I\pm J_R)/2$,
\[
 \boxed{\mathcal T_{p,a}(\sigma)=
 \|P_+v_{p,a,\sigma}\|^2-\|P_-v_{p,a,\sigma}\|^2.}            \tag{2.10}\label{eq:parity}
\]

The map $e_t\mapsto e_{t/R}$ embeds this shell space into the finitely supported
part of $\ell^2(\mathbb Q/\mathbb Z)$.  Its support consists of points of exact
additive order $R$, so supports belonging to distinct residuals are disjoint.
For
\[
 V_{p,\sigma}=\sum_{p/4<a<p/2}v_{p,a,\sigma}
\]
one therefore has
\[
 \langle V_{p,\sigma},JV_{p,\sigma}\rangle
 =\sum_{p/4<a<p/2}\mathcal T_{p,a}(\sigma).                       \tag{2.11}
\]
On finitely supported functions define $(A_qf)(x)=f(qx)$.  Then
\[
 \boxed{A_q^*JA_q=qJ.}                                           \tag{2.12}\label{eq:degree}
\]
In particular the degree action preserves signed isotropy.
\end{proposition}

\begin{proof}
Expand $|Z_{S,\chi}|^2$ with the real positive divisor weights.  Character
orthogonality applied to $-de^{-1}\in G_R$ retains exactly $d+e\equiv0\pmod R$,
which proves \eqref{eq:char}.  The self-adjoint involution $J_R$ has spectral
projections $P_\pm$, giving \eqref{eq:parity}.  Since $R$ is odd, both
$e_t+e_{-t}$ and $e_t-e_{-t}$ occur, so the signed form is indefinite; it cannot
be replaced on the full space by the pullback of a positive semidefinite form.

For $t\in G_R$, the class $t/R$ has exact order $R$.  This proves the disjoint
support assertion and (2.11).  Finally, each $y\in\mathbb Q/\mathbb Z$ has
exactly $q$ preimages under multiplication by $q$.  Hence
$\langle A_qf,A_qg\rangle=q\langle f,g\rangle$, while $A_qJ=JA_q$.
Equation \eqref{eq:degree} follows.
\end{proof}

The operator in (2.12) is the unscaled global pullback; explicitly,
\[
 (A_q^*g)(y)=\sum_{qx=y}g(x),\qquad A_q^*A_q=qI.                 \tag{2.13}
\]
The scaled operator $V_q=q^{-1/2}A_q$ instead satisfies $V_q^*JV_q=J$.  The global
operator need not preserve one exact-order stratum: its compression to the
order-$R$ stratum is the permutation $t\mapsto qt$ when $\gcd(q,R)=1$, and is
zero when $\gcd(q,R)>1$; all uncompressed support remains in its actual
exact-order strata.

Equations \eqref{eq:char}--\eqref{eq:degree} locate the sign that any analytic
comparison must retain.  A lower bound for the unsigned energy does not imply a
strict lower bound for \eqref{eq:parity}.

\section{Cross-shell boundary and its obstruction space}

Choose an integer $m\ge1$, put $q=4m+1$, and assume $qR<p$.  Retain two original
shells
\[
 a=\frac{p+R}{4},\qquad a'=\frac{p+qR}{4}=a+mR.                  \tag{3.1}
\]
Take actual positive divisors
\[
 d\mid pa,\qquad e\mid pa',\qquad R\mid d+e,                   \tag{3.2}
\]
and their actual complements
\[
 \alpha=\frac{pa}{d},\qquad \beta=\frac{pa'}e.                  \tag{3.3}
\]
By (1.1), $d$ is a unit modulo $R$.  Since
$e\beta-d\alpha=pa'-pa=pmR$ and $e\equiv-d\pmod R$, cancellation
of $d$ gives $R\mid\alpha+\beta$.  Thus the following are positive integers:
\[
 \lambda=\frac{d+e}{R},\qquad \mu=\frac{\alpha+\beta}{R}.       \tag{3.4}
\]

\begin{theorem}[Exact cross-shell boundary]\label{thm:cross-shell}
The retained data satisfy both determinant identities
\[
 \boxed{\lambda\beta-d\mu=pm,\qquad e\mu-\lambda\alpha=pm}     \tag{3.5}
\]
and the four-term reciprocal identity
\[
 \boxed{
 \frac4p=\frac1a+\frac1{\lambda\alpha}+\frac1{\mu d}
 +\frac{m}{a\lambda\mu}.}                                      \tag{3.6}\label{eq:fourterm}
\]
Let
\[
 L=\lambda\mu,\quad g=\gcd(L,m),\quad
 \ell=\frac Lg,\quad \nu=\frac{L+m}{g}.                        \tag{3.7}
\]
Then $\gcd(\ell,\nu)=1$, and the first and fourth terms in
\eqref{eq:fourterm} form one unit fraction exactly when
\[
 \boxed{\nu\mid a.}                                             \tag{3.8}\label{eq:fusion}
\]
On this exact domain,
\[
 \boxed{
 X=\frac{a\lambda\mu}{\lambda\mu+m}
   =\frac a\nu\frac{\lambda\mu}{g},\quad
 Y=\lambda\alpha,\quad Z=\mu d}                                \tag{3.9}
\]
are positive integers and $4/p=1/X+1/Y+1/Z$.
\end{theorem}

\begin{proof}
The first identity in (3.5) is
\[
 \lambda\beta-d\mu
 =\frac{(d+e)\beta-d(\alpha+\beta)}R
 =\frac{e\beta-d\alpha}{R}=pm.
\]
The second is obtained by retaining the other column.  Since $d\alpha=pa$,
\[
 \frac1{\lambda\alpha}+\frac1{\mu d}
 +\frac{m}{a\lambda\mu}
 =\frac{\mu d+\lambda\alpha+pm}{pa\lambda\mu}
 =\frac{\lambda(\alpha+\beta)}{pa\lambda\mu}
 =\frac R{pa}.
\]
Adding $1/a$ and retaining $p+R=4a$ proves \eqref{eq:fourterm}.

The retained fusion is
\[
 \frac1a+\frac{m}{aL}=\frac{L+m}{aL}=\frac{\nu}{a\ell}.
\]
Because $\gcd(\nu,\ell)=1$, its reciprocal $a\ell/\nu$ is integral exactly
when $\nu\mid a$.  Equations (3.9) then retain all three denominators.
\end{proof}

\begin{definition}[Boundary obstruction]\label{def:boundary-obstruction}
For fixed $m,\lambda,\mu$, let
\[
 \mathcal O_\nu=(\nu^{-1}\mathbb Z)/\mathbb Z,\qquad
 \Theta_{\ell,\nu}:\mathbb Z\longrightarrow\mathcal O_\nu,\quad
 n\longmapsto\frac{n\ell}{\nu}+\mathbb Z.                       \tag{3.10}
\]
The exact sequence
\[
 0\longrightarrow\nu\mathbb Z\longrightarrow\mathbb Z
 \xrightarrow{\Theta_{\ell,\nu}}\mathcal O_\nu\longrightarrow0 \tag{3.11}
\]
defines the cyclic obstruction class
\[
 \Omega(a;m,\lambda,\mu)
 =\Theta_{\ell,\nu}(a)
 =\frac{a\lambda\mu}{\lambda\mu+m}+\mathbb Z.                  \tag{3.12}
\]
It vanishes exactly on the fusion domain \eqref{eq:fusion}.
\end{definition}

Put $h=\gcd(a,\nu)$, $A=a/h$, and $\rho=\nu/h$.  The complete retained defect is
\[
 \frac1a+\frac{m}{aL}=\frac{\rho}{A\ell},\qquad
 \gcd(\rho,A\ell)=1,
 \qquad \ord\Omega=\rho.                                       \tag{3.13}
\]
Thus a failed fusion is not discarded: its exact reduced numerator is the
order of a point in $\mathcal O_\nu$.  The order-$\rho$ stratum is
\[
 \mathcal S_\rho=
 \left\{\frac{\nu}{\rho}w\pmod\nu:
 w\in(\mathbb Z/\rho\mathbb Z)^\times\right\},                  \tag{3.14}
\]
with $\varphi(\rho)$ residue classes.  In the original shell variables,
\[
 \nu\mid a\Longleftrightarrow4\nu\mid p+R,\qquad
 \rho=\frac{4\nu}{\gcd(p+R,4\nu)}.                             \tag{3.15}
\]
The surviving identity outside the fusion locus remains
\[
 \frac4p=\frac1{\lambda\alpha}+\frac1{\mu d}
 +\frac{\rho}{A\ell}.                                           \tag{3.16}
\]

\begin{figure}[ht]
\centering
\begin{tikzpicture}[>=Latex,node distance=17mm,
  box/.style={draw,rounded corners,align=center,inner sep=4pt},
  map/.style={->,thick}]
\node[box] (left) {shell $a$\\$d\mid pa$, $\alpha=pa/d$};
\node[box,right=28mm of left] (right) {shell $a'=a+mR$\\$e\mid pa'$, $\beta=pa'/e$};
\node[box,below=18mm of $(left)!0.5!(right)$] (bridge)
 {$\lambda=(d+e)/R$, $\mu=(\alpha+\beta)/R$\\
 $\lambda\beta-d\mu=pm$};
\node[box,below left=19mm and 11mm of bridge] (four)
 {$\dfrac4p=\dfrac1a+\dfrac1{\lambda\alpha}
 +\dfrac1{\mu d}+\dfrac{m}{a\lambda\mu}$};
\node[box,below right=19mm and 11mm of bridge] (obs)
 {$\Omega=\dfrac{a\lambda\mu}{\lambda\mu+m}+\mathbb Z$\\
 $\ord(\Omega)=\rho$};
\node[box,below=19mm of $(four)!0.5!(obs)$] (return)
 {$\Omega=0\iff\nu\mid a$\\
 $(X,Y,Z)=\left(\dfrac{a\lambda\mu}{\lambda\mu+m},
 \lambda\alpha,\mu d\right)$};
\draw[map] (left) -- node[above] {$R\mid d+e$} (right);
\draw[map] (left) -- (bridge);
\draw[map] (right) -- (bridge);
\draw[map] (bridge) -- (four);
\draw[map] (bridge) -- (obs);
\draw[map] (four) -- (return);
\draw[map] (obs) -- (return);
\end{tikzpicture}

\caption{Every displayed arrow retains its domain and codomain.  The cross-shell
pair produces the determinant $pm$ and the fourth reciprocal term.  Fusion is
an exact map only on the kernel of $\Theta_{\ell,\nu}$; outside it, the order
$\rho$ records the surviving boundary numerator.}
\label{fig:cross-shell}
\end{figure}

\section{General odd-quotient complement law}

The same-shell theorem is best understood before setting the quotient equal to
three.  Continue with one original shell $(p,a,R)$.  Suppose a positive divisor
$m\mid R$ satisfies
\[
 \gcd(R,p-1)=m,\qquad R\mid m^2,\qquad N=R/m>1.               \tag{4.1}\label{eq:Ndomain}
\]
All of $R,m,N$ are odd.  Define the complete trace source
\[
 \mathcal T_N=\{u\mid a^2:m\mid p+4u\},                         \tag{4.2}
\]
the complement $\iota(u)=a^2/u$, and
\[
 b_u=\frac{p+4u}{m}\pmod N,\qquad
 \omega=\frac{p-1}{m}\pmod N.                                 \tag{4.3}
\]

\begin{theorem}[Exact $N$-colour complement decomposition]
\label{thm:N-colour}
Under \eqref{eq:Ndomain}, $\omega$ is a unit in $\mathbb Z/N\mathbb Z$,
$\iota$ is a fixed-point-free involution of $\mathcal T_N$, and
\[
 \boxed{b_{\iota(u)}=-b_u\pmod N.}                              \tag{4.4}\label{eq:complement-colour}
\]
The full gates are
\[
 u\in M_a\Longleftrightarrow b_u=0,\qquad
 u\in E_a\Longleftrightarrow b_u=\omega.                        \tag{4.5}
\]
Let
\[
 U_N=\{u\in\mathcal T_N:b_u\notin\{0,\omega,-\omega\}\}.
\]
Then $U_N$ is complement-stable, and there is an exact labelled orbit
decomposition
\[
 \boxed{
 \mathcal T_N/\iota\ \simeq\
 E_a\ \amalg\ (M_a/\iota)\ \amalg\ (U_N/\iota).}              \tag{4.6}\label{eq:orbit-decomp}
\]
Consequently
\[
 \boxed{|\mathcal T_N|=2|E_a|+|M_a|+|U_N|.}                     \tag{4.7}\label{eq:defect-count}
\]
The inverse in \eqref{eq:orbit-decomp} attaches the actual pair
$\{u,a^2/u\}$; an input-member bit restores the labelled source word.
\end{theorem}

\begin{proof}
The exact gcd in \eqref{eq:Ndomain} gives
\[
 \gcd\left(N,\frac{p-1}{m}\right)=1,
\]
so $\omega$ is a unit.  For $u\in\mathcal T_N$,
\[
 4(a+u)=p+R+4u
\]
is divisible by $m$.  Since $m$ is odd, $m\mid a+u$; since $R\mid m^2$,
$R\mid(a+u)^2$.  The divisor $u$ is a unit modulo $R$, and
\[
 u+\frac{a^2}{u}+2a=\frac{(a+u)^2}{u}\equiv0\pmod R.           \tag{4.8}
\]
Therefore
\[
 (p+4u)+\left(p+4\frac{a^2}{u}\right)\equiv0\pmod R,
\]
which proves that $\iota(u)\in\mathcal T_N$ and gives
\eqref{eq:complement-colour} after division by $m$.

A fixed point would have $u=a$.  Membership in $\mathcal T_N$ would then give
$m\mid p+4a=2p+R$, hence $m\mid2p$.  But $m>1$ is odd and is coprime to $p$,
a contradiction.  Finally,
\[
 p+4u-(p-1)=4u+1
\]
proves (4.5).  Every orbit with colour zero consists of two middle words; every
orbit with colours $\{\omega,-\omega\}$ contains exactly one exterior word;
the remaining orbits are exactly $U_N/\iota$.  This proves (4.6)--(4.7).
\end{proof}

\begin{corollary}[Uniqueness of three colours]\label{cor:unique-three}
Complementation alone guarantees that every possible colour orbit meets a full
E or M gate if and only if $N=3$.  More exactly,
\[
 \{0,\omega,-\omega\}=\mathbb Z/N\mathbb Z
 \Longleftrightarrow N=3.                                      \tag{4.9}
\]
For a particular source at $N>3$, return still holds exactly when its realized
colour support is contained in $\{0,\omega,-\omega\}$; the nonempty set $U_N$
is the exact obstruction otherwise.
\end{corollary}

\begin{proof}
The three displayed residues are distinct because $N>1$ is odd and $\omega$ is
a unit.  Equality with the entire residue group therefore forces $N=3$.
Conversely a group of order three consists exactly of those residues.
\end{proof}

\begin{proposition}[Sharpness for every larger odd quotient]
\label{prop:larger-quotient-sharp}
For every odd $N>3$, there are infinitely many primes $p\equiv1\pmod{24}$
and original shells satisfying \eqref{eq:Ndomain} whose source
$\mathcal T_N$ contains a wholly uncovered complement orbit.
\end{proposition}

\begin{proof}
Set
\[
 m=3N,\qquad R=3N^2=mN.                                        \tag{4.14}
\]
Then $R\equiv3\pmod4$ and $R\mid m^2$.  By the Chinese remainder theorem,
choose $s>0$ with
\[
 s\equiv1\pmod N,\qquad ms\equiv5\pmod8,                      \tag{4.15}
\]
and retain
\[
 U=\frac{ms-1}{4}.                                               \tag{4.16}
\]
Thus $U$ is odd and $\gcd(U,mN)=1$.  Choose an even $k_0$ with
\[
 4Uk_0\equiv2\pmod N,                                          \tag{4.17}
\]
which is possible because $4U$ is a unit modulo odd $N$ and adding $N$ changes
parity.  For $k=k_0+2Nj$, define
\[
 a=U(mk-1),\qquad p=4a-R.                                      \tag{4.18}
\]
As $j$ varies, $p$ lies in the progression
\[
 p=p_0+8UmN\,j.                                                  \tag{4.19}
\]
The retained identities give $p\equiv1\pmod8$ and $p\equiv1\pmod3$, hence
$p\equiv1\pmod{24}$.  Moreover $p\equiv1\pmod m$,
$p\equiv-R\pmod U$, and $\gcd(U,R)=1$, so
\[
 \gcd(p_0,8UmN)=1.
\]
Dirichlet's theorem on primes in reduced progressions
\citep[Theorem~18.1]{Sutherland} supplies infinitely many prime values; choose
one with $p>R$.

Using $4U=ms-1$, one has
\[
 p-1=m(4Uk-s-N),\qquad \frac{p-1}{m}\equiv1\pmod N.             \tag{4.20}
\]
Therefore $\gcd(R,p-1)=m$ and $\omega=1$.  The actual divisor $U\mid a^2$
satisfies
\[
 p+4U=m(4Uk-N),\qquad b_U=2\pmod N.                              \tag{4.21}
\]
Its complement has colour $-2$.  For odd $N>3$, neither colour is $0$ or
$\omega=1$, so this entire complement orbit lies in $U_N$.
\end{proof}

\begin{remark}[The quotient-one edge]
If $N=1$, then $m=R$ and $\mathcal T_1=M_a=E_a$: the two gates coincide because
$R\mid p-1$.  In the original shell domain $R\equiv3\pmod4$, hence $R>1$, and
the divisor complement is free.  This doubly gated edge is separate from the
nontrivial odd-quotient uniqueness in Corollary~\ref{cor:unique-three}.
\end{remark}

\subsection{Exact morphism to the affine channel coordinate}

Let
\[
 K=\prod_{\ell^e\parallel R}\ell^{\lceil e/2\rceil},\qquad
 D=R/K.                                                         \tag{4.10}
\]
The hypotheses $R=mN$ and $R\mid m^2$ first give $N\mid m$.  They also imply,
prime by prime,
\[
 K\mid m,\qquad N\mid K.                                      \tag{4.11}
\]
Write $t=m/K$; then $D=Nt$.  The source $\mathcal T_N$ lies in the common E/M
trace locus of the affine-channel theorem because $K\mid m$ divides both
$p-1$ and $p+4u$.  Its middle fine coordinate is
\[
 c_M\equiv p^{-1}\frac{p+4u}{K}\pmod D.
\]
Since $D\mid m$ and $m\mid p-1$, one has $p^{-1}\equiv1\pmod D$.  Thus
\[
 c_M=t\,b_u\pmod{Nt}.                                          \tag{4.12}
\]
The exact subgroup quotient
\[
 t\mathbb Z/(Nt)\mathbb Z\longrightarrow\mathbb Z/N\mathbb Z,
 \qquad tb\longmapsto b                                       \tag{4.13}
\]
sends the prior affine E target
$(p-1)/K=t\omega$ to $\omega$, the M target to zero, and the divisor complement
to $b\mapsto-b$.  This is the exact morphism between the fine-defect channel
operator and the colour-orbit theorem; neither presentation replaces the other.

\begin{figure}[ht]
\centering
\begin{tikzpicture}[>=Latex,node distance=17mm,
  colour/.style={circle,draw,minimum size=10mm,inner sep=1pt},
  gate/.style={rounded corners,draw,align=center,inner sep=3pt}]
\node[colour,fill=blue!10] (zero) {$0$};
\node[colour,fill=green!12,right=31mm of zero] (omega) {$\omega$};
\node[colour,fill=orange!15,right=31mm of omega] (minus) {$-\omega$};
\draw[->,bend left=22] (omega) to node[above] {$u\mapsto a^2/u$} (minus);
\draw[->,bend left=22] (minus) to node[below] {$u\mapsto a^2/u$} (omega);
\draw[->,loop above] (zero) to node[above] {colour fixed} (zero);
\node[gate,below=14mm of zero] (M) {M gate\\two divisor words};
\node[gate,below=14mm of omega] (E) {E gate\\one divisor word};
\node[gate,below=14mm of minus] (Ec) {complement\\of the E word};
\draw[->] (zero) -- (M);
\draw[->] (omega) -- (E);
\draw[->] (minus) -- (Ec);
\node[draw,dashed,fit=(zero)(omega)(minus),inner sep=5mm,
 label=above:{${\mathbb Z}/3{\mathbb Z}=\{0,\omega,-\omega\}$}] {};
\end{tikzpicture}

\caption{The exact quotient-three mechanism.  At residue level, complementation
fixes colour zero and exchanges $\omega$ with $-\omega$; on divisor words it is
fixed-point-free.  The zero orbit contains two M words, while the nonzero orbit
contains exactly one E word.  Larger odd quotients have additional uncovered
orbits.}
\label{fig:three-colour}
\end{figure}

\section{Same-shell third-defect return}

Assume now
\[
 \boxed{\gcd(R,p-1)=R/3},\qquad m=R/3.                           \tag{5.1}
\]
Because $p\equiv1\pmod3$, valuation comparison forces $9\mid R$; hence
$3\mid m$ and $R=3m\mid m^2$.  Thus Theorem~\ref{thm:N-colour} applies with
$N=3$ to the complete source
\[
 \mathcal T_3=\{u\mid a^2:m\mid p+4u\}.                         \tag{5.2}
\]

\begin{theorem}[Whole-shell third-defect completion]\label{thm:third-defect}
Every member of $\mathcal T_3$ gives a rational middle trace whose two tails
have denominators dividing three.  It returns at unchanged $(p,R,a)$ to an
original integer E or M word using either $u$ or $a^2/u$.  Explicitly, with
\[
 b_u=\frac{p+4u}{m}\pmod3,\qquad
 \omega=\frac{p-1}{m}\pmod3,
\]
the return is
\[
 \begin{array}{c|c}
 b_u&\text{marked integer output}\\ \hline
 0&M\text{ at }u\\
 \omega&E\text{ at }u\\
 -\omega&E\text{ at }a^2/u.
 \end{array}                                                     \tag{5.3}
\]
Moreover
\[
 \boxed{2|E_a|+|M_a|=|\mathcal T_3|},                           \tag{5.4}
\]
and therefore
\[
 \boxed{E_a\cup M_a=\varnothing\Longleftrightarrow
 \mathcal T_3=\varnothing.}                                    \tag{5.5}
\]
The number of increasing original denominator triples at this shell is exactly
$|\mathcal T_3|/2$.
\end{theorem}

\begin{proof}
Only the tail and increasing-order assertions remain beyond
Theorem~\ref{thm:N-colour} and Corollary~\ref{cor:unique-three}.  Put
$v=a^2/u$.  The rational middle tails are retained as
\[
 y=\frac{p(a+v)}R,\qquad z=\frac{p(a+u)}R.                       \tag{5.6}
\]
Both $u$ and $v$ belong to $\mathcal T_3$, so $m$ divides both $a+u$ and
$a+v$; since $R=3m$, each tail has reduced denominator dividing three.  Their
sum retains the exact integer formula
\[
 y+z=\frac{p(a+u)^2}{Ru}.                                       \tag{5.7}
\]
Indeed, $R\mid(a+u)^2$, $u\mid(a+u)^2$, and $\gcd(R,u)=1$.

For an exterior word, (1.5) gives $a<y<z$: the inequalities are equivalent to
$\kappa>r$ and $pr>s$, which follow from $p>R$ and $s<a<p$.  Exactly one tail
is divisible by $p$: the factor $p$ is explicit in $z$, whereas
$R\kappa=pr+s$ and $0<s<p$ show that $p\nmid\kappa$, hence $p\nmid y$.
The ordered inverse recovers its unique exterior word.
For a middle word, both tails in (1.6) exceed $a$ and are divisible by $p$;
complementation swaps $r,s$ and hence swaps the two tails.  There is no fixed
middle word in $\mathcal T_3$.  Therefore exterior words count once and middle
complement pairs count once:
\[
 \#\{\text{increasing triples}\}=|E_a|+\frac{|M_a|}{2}
 =\frac{|\mathcal T_3|}{2}.
\]
\end{proof}

Under complementation, the primitive coordinates change exactly by
\[
 (h,r,s)\longmapsto(h,s,r).                                     \tag{5.8}
\]
The prime, residual, distinguished denominator, and original exponent box are
unchanged; the E/M quotient is recomputed in its own channel.

\section{All applicable residuals from the prescribed prime}

Write the complete factorization
\[
 p-1=2^k3^\eta b,\qquad \gcd(b,6)=1,\qquad k\ge3,\quad\eta\ge1. \tag{6.1}
\]

\begin{theorem}[Applicable residual classification]\label{thm:residuals}
The positive residuals $R<p$, $R\equiv3\pmod4$ satisfying
$\gcd(R,p-1)=R/3$ are exactly
\[
 \boxed{R=3^{\eta+1}d,\qquad d\mid b,\qquad
 d\equiv(-1)^\eta\pmod4.}                                     \tag{6.2}
\]
Every one satisfies
\[
 R\le\frac{3(p-1)}{2^k}<p.                                    \tag{6.3}
\]
The family is nonempty exactly when $\eta$ is even or $b$ has a prime divisor
congruent to $3$ modulo $4$.
\end{theorem}

\begin{proof}
At the prime three, the equality
\[
 \min(v_3(R),\eta)=v_3(R)-1
\]
forces $v_3(R)=\eta+1$.  At every other odd prime, its exponent in $R$ is at
most its exponent in $p-1$.  This yields $R=3^{\eta+1}d$ with $d\mid b$.
The retained shell congruence $R\equiv3\pmod4$ is exactly
$d\equiv(-1)^\eta\pmod4$.  These conditions also prove the converse gcd
identity.  Since $d\le b$, (6.3) follows.  If $\eta$ is even, $d=1$ works.  If
$\eta$ is odd, a divisor $d\equiv3\pmod4$ exists exactly when one of the prime
divisors of $b$ is $3$ modulo $4$.
\end{proof}

This theorem constructs every candidate residual from the prescribed prime.
It does not assert that one of the associated sets $\mathcal T_3$ is nonempty.

\section{Complete factor-level failure locus at residual 27}

\begin{theorem}[Residual-27 classification]\label{thm:R27}
Let $p$ be prime with
\[
 \boxed{p\equiv73\ \text{or}\ 145\pmod{216}},\qquad
 a=\frac{p+27}{4}.                                               \tag{7.1}
\]
The complete original $R=27$ shell is empty if and only if one of the following
disjoint factorization types occurs:
\begin{enumerate}
\item every prime factor of $a$ is $1$ modulo $3$;
\item exactly two prime-factor occurrences of $a$ are $5$ modulo $9$, counting
multiplicity, and every remaining prime factor is $1$ modulo $9$.
\end{enumerate}
Outside these two types, an original E or M word is constructed with at most
three signed prime occurrences in $u/a$, followed by (5.3).
\end{theorem}

\begin{proof}
The two classes in (7.1) are exactly the classes $p\equiv1\pmod{24}$ for which
$\gcd(27,p-1)=9$.  They give $a\equiv7\pmod9$ and
\[
 \mathcal T_3=\{u\mid a^2:u\equiv2\pmod9\}.                    \tag{7.2}
\]
Write
\[
 a=\prod_q q^{e_q},\qquad
 u=\prod_q q^{e_q+\beta_q},\qquad -e_q\le\beta_q\le e_q.      \tag{7.3}
\]
The exact target is
\[
 \prod_q q^{\beta_q}\equiv-1\pmod9.                            \tag{7.4}
\]
The unit group modulo nine is generated by two; the residues
\[
 1,2,4,8,7,5
\]
have logarithms $0,1,2,3,4,5$ modulo six.  The logarithm of $a$ is four and
the target in (7.4) is three.

A residue-$8$ prime supplies the target with one occurrence.  Suppose none is
present.  Let $H$ be the total multiplicity of primes in the residue classes
$2,5$ modulo nine.  The logarithm of $a$ is even, so $H$ is even.  If $H\ge4$,
select three such occurrences, assigning sign $+1$ to residue two and sign
$-1$ to residue five; each contributes logarithm one.  If $H=2$ and a residue
$4$ or $7$ occurrence exists, one signed odd logarithm contributes one and one
signed even logarithm contributes two.

If $H=0$, every available logarithm is even, so the target three is impossible;
this is the first exceptional type.  In the remaining case $H=2$, with no
residue $8,4,$ or $7$, the total logarithm four forces both odd occurrences to
be residue five.  Their complete centred exponent range contributes only
\[
 \{-2,-1,0,1,2\}\pmod6,
\]
which misses three.  This is the second type.  Every successful selection obeys
the literal bounds in (7.3), so it returns an actual divisor of $a^2$.
\end{proof}

The exceptional examples
\[
 (p,a)=(1009,259=7\cdot37),\qquad (p,a)=(73,25=5^2)              \tag{7.5}
\]
have, respectively, the first and second factorization type.  Their residual-27
shells are empty; neither prime is an Erd\H{o}s--Straus counterexample.

\section{What remains}

The results above are structural rather than a finite prime census.  They give:
\begin{enumerate}
\item a complete source that exists at every shell and whose signed value is
strictly positive exactly when that shell has an original integer word;
\item an exact cross-shell map with a classified cyclic boundary obstruction;
\item an exact same-shell quotient law in which three is the unique nontrivial
odd quotient closed by the divisor complement alone;
\item every residual on which that three-colour theorem applies, constructed
from the prescribed prime; and
\item the complete emptiness criterion for the first canonical residual $27$.
\end{enumerate}
They do not prove that an applicable $\mathcal T_3$ is nonempty for every hard
prime, nor that a cross-shell pair with zero boundary obstruction must occur.
An arbitrary counterexample would have to make every constructed three-colour
source empty, defeat all other original shells, and avoid every cross-shell
kernel in (3.11).  Proving that this simultaneous avoidance is impossible is
the remaining existence problem.

\section*{Reproducibility and attribution}

The received four-page proof and constructor are preserved byte-for-byte under
\texttt{received/}.  The active constructor corrects two output-field names but
does not alter the mathematics.  A separately written checker reconstructs the
complement law, count, residual family, residual-27 factor classification,
signed pairing at fixed exact test values, and cross-shell identities without
importing the active constructor.  The Boolean-pair antecedent is credited to
the JT collaboration; the type-I/type-II antecedent is credited to Elsholtz and
Tao.  No historical-priority claim is made for the new quotient and obstruction
formulations.
